\section{Thiết kế và đánh giá sản phẩm Demo}

\subsection{Kết quả khảo sát nhu cầu thực tế của giảng viên tại Đại học Bách khoa Hà Nội}
Để xây dựng một giải pháp công nghệ đáp ứng đúng yêu cầu thực tiễn, nghiên cứu đã tiến hành khảo sát nhu cầu thực tế của các giảng viên đang tham gia giảng dạy tại Đại học Bách khoa Hà Nội. Nội dung khảo sát tập trung vào việc đánh giá quỹ thời gian dành cho công tác chuẩn bị bài giảng, xác định các điểm nghẽn tiêu tốn nhiều công sức nhất và đo lường mức độ sẵn sàng đón nhận công nghệ trí tuệ nhân tạo tạo sinh trong hoạt động học thuật.

Kết quả thu thập được cho thấy đa số giảng viên phải đối mặt với áp lực thời gian lớn khi thực hiện công tác chuẩn bị học liệu hàng tuần. Cụ thể, một nửa số giảng viên tham gia khảo sát cho biết họ phải dành từ năm đến mười giờ mỗi tuần để nghiên cứu và thiết kế nội dung bài giảng. Bên cạnh đó, có một tỷ lệ giảng viên phải tiêu tốn trên mười giờ mỗi tuần cho công tác này. Số lượng giảng viên có thể hoàn thành việc soạn thảo học liệu dưới năm giờ chiếm tỷ lệ ít hơn. Điều này chứng minh rằng việc chuẩn bị bài giảng cho các học phần chuyên sâu thuộc khối kỹ thuật và công nghệ yêu cầu một quỹ thời gian lớn từ người dạy.

Khi đi sâu phân tích các công đoạn cụ thể trong quy trình soạn bài, kết quả khảo sát xác định hai công đoạn tiêu tốn nhiều thời gian và công sức nhất của giảng viên là đọc, chắt lọc, tóm tắt lượng lớn tài liệu chuyên môn và cấu trúc lại kiến thức thành các ý chính để đưa lên bản trình chiếu. Các công đoạn khác như tìm kiếm ví dụ thực tế hoặc rà soát cập nhật thông tin tuy có đòi hỏi công sức nhưng không phải là nguyên nhân chính gây ra áp lực thời gian. Áp lực từ việc tổng hợp tài liệu chuyên môn sâu cũng được phản ánh qua nhận định của giảng viên, khi có đến hơn tám mươi phần trăm giảng viên hoàn toàn đồng ý hoặc đồng ý rằng hoạt động này tạo ra áp lực lớn về mặt thời gian cho họ.

Về mức độ sẵn sàng ứng dụng nội dung do trí tuệ nhân tạo tạo ra, phần lớn giảng viên thể hiện thái độ cẩn trọng. Khoảng hai phần ba số giảng viên được hỏi cho biết họ khá sẵn sàng sử dụng các nội dung do công nghệ tạo ra nhưng chỉ dưới dạng một bản nháp thô để tự tinh chỉnh, thay vì đưa vào sử dụng trực tiếp mà không qua kiểm duyệt. Xu hướng này cho thấy giảng viên ý thức rõ tầm quan trọng của việc kiểm soát chất lượng thông tin học thuật và mong muốn giữ vai trò chủ đạo trong việc quyết định nội dung bài giảng.

Cuối cùng, khi được hỏi về các tính năng ưu tiên phát triển cho một sản phẩm hỗ trợ học liệu, giảng viên tập trung lựa chọn hai tính năng cốt lõi là tóm tắt tài liệu dài thành dàn ý bài giảng và chuyển đổi văn bản thành cấu trúc ý cô đọng làm bản trình chiếu. Tính năng tự động sinh ngân hàng câu hỏi đánh giá đa dạng cũng được quan tâm nhưng xếp sau hai nhu cầu trên. Kết quả này là cơ sở thực hiện thiết kế cấu trúc chức năng và các nguyên tắc xây dựng cho sản phẩm thử nghiệm.

\subsection{Các nguyên tắc trong xây dựng sản phẩm thử nghiệm}
Dựa trên kết quả khảo sát thực tế và các khung lý thuyết sư phạm, nghiên cứu đề xuất bốn nguyên tắc cốt lõi làm nền tảng cho việc thiết kế và triển khai công cụ thử nghiệm hỗ trợ giảng viên:

Nguyên tắc thứ nhất là lấy người dạy làm trung tâm, đặt giảng viên vào vị trí kiểm duyệt trực tiếp toàn bộ quy trình. Công cụ trí tuệ nhân tạo tạo sinh đóng vai trò một trợ lý xây dựng bản thảo và gợi ý ý tưởng ban đầu. Giảng viên giữ quyền quyết định cuối cùng, thực hiện rà soát, hiệu chỉnh và phê duyệt mọi nội dung trước khi đưa vào giảng dạy. Nguyên tắc này giúp giải quyết lo ngại của giảng viên về tính chính xác của dữ liệu và bảo đảm vai trò định hướng chuyên môn của người dạy trong môi trường sư phạm.

Nguyên tắc thứ hai là chuẩn hóa sư phạm. Sản phẩm không đơn thuần là một công cụ xử lý văn bản tự động mà được thiết kế tuân thủ các chuẩn mực sư phạm. Cụ thể, chức năng sinh câu hỏi kiểm tra được ràng buộc tương thích với các bậc nhận thức của thang đo Bloom cải tiến. Đồng thời, các chức năng gợi ý hoạt động dạy học bám sát thuyết kiến tạo, tập trung vào việc kích thích sự tương tác của sinh viên và thúc đẩy tư duy giải quyết vấn đề thực tế.

Nguyên tắc thứ ba là bảo đảm liêm chính học thuật và an toàn thông tin. Việc sử dụng các mô hình ngôn ngữ lớn để xử lý tài liệu giảng dạy đòi hỏi cơ chế bảo mật dữ liệu nhằm bảo vệ quyền sở hữu trí tuệ của giảng viên đối với các bài giảng chuyên ngành. Đồng thời, thiết kế kỹ thuật của hệ thống áp dụng các chuỗi câu lệnh gợi ý chặt chẽ để hạn chế hiện tượng ảo tưởng thông tin của mô hình ngôn ngữ, bảo đảm kiến thức kỹ thuật được cung cấp cho sinh viên Đại học Bách khoa Hà Nội luôn chính xác.

Nguyên tắc thứ tư là tích hợp tự động hóa toàn diện từ một nguồn dữ liệu duy nhất. Thay vì yêu cầu giảng viên sử dụng rời rạc nhiều công cụ khác nhau cho từng tác vụ soạn bài, hệ thống được thiết kế để tiếp nhận một tài liệu học thuật nguồn duy nhất và tự động chuyển hóa thành hệ thống học liệu đa dạng, đồng bộ từ đề cương, slide bài giảng, ngân hàng câu hỏi cho đến kịch bản tương tác trên lớp.

\subsection{Ứng dụng lý thuyết vào thiết kế và triển khai sản phẩm thực tế ``Không gian Học liệu số HUST''}

\subsubsection{Tổng quan kiến trúc hệ thống và quy trình xử lý toàn diện}
Từ các nguyên tắc thiết kế và kết quả khảo sát thực tế, nghiên cứu đã tiến hành phát triển sản phẩm web thử nghiệm với tên gọi \textbf{``Không gian Học liệu số HUST''}. Hệ thống được xây dựng nhằm cung cấp giải pháp tự động khởi tạo học liệu toàn diện dành cho giảng viên Đại học Bách khoa Hà Nội.

Khi giảng viên tải lên một tài liệu học thuật nguồn (như giáo trình, bài báo khoa học hay ghi chú bài giảng), hệ thống thực hiện quy trình trích xuất và phân rã ngữ cảnh. Thông qua các chuỗi câu lệnh gợi ý được thiết kế tối ưu, hệ thống hạn chế tối đa sai lệch thông tin và tự động khởi tạo bốn phân hệ học liệu đồng bộ.

\subsubsection{Triển khai bốn phân hệ học liệu cốt lõi}

\subsubsection*{Phân hệ 1: Bản đồ tri thức và đề cương bài giảng dưới dạng Sơ đồ tư duy tương tác}
Phân hệ đầu tiên chuyển đổi văn bản thô của tài liệu nguồn thành một sơ đồ tư duy tri thức tương tác. Cấu trúc bài giảng được tổ chức phân cấp từ các chủ đề chính đến các khái niệm thành phần. Điểm đặc biệt của phân hệ này là tính năng tương tác đa chiều: giảng viên và sinh viên có thể nhấp chuột vào từng nút tri thức trên sơ đồ để mở rộng cửa sổ hiển thị nội dung chi tiết. Đồng thời, hệ thống tích hợp không gian thảo luận trực tiếp ngay tại từng nút tri thức. Thiết kế này hiện thực hóa nguyên lý của thuyết kiến tạo, cho phép người học chủ động đào sâu từng mảng kiến thức và trao đổi ý kiến chuyên môn trực tiếp trên cấu trúc sơ đồ bài học.

\subsubsection*{Phân hệ 2: Khởi tạo kịch bản Slide bài giảng}
Phân hệ thứ hai giải quyết điểm nghẽn tiêu tốn nhiều thời gian nhất của giảng viên theo khảo sát: công đoạn cấu trúc lại kiến thức để đưa lên bản trình chiếu. Dựa trên nội dung tài liệu nguồn, hệ thống tự động chắt lọc các thuật ngữ quan trọng, phân rã các đoạn văn dài thành các cấu trúc ý ngắn gọn, súc tích và sắp xếp theo trình tự kịch bản giảng dạy. Sản phẩm đầu ra của phân hệ này giúp giảng viên nhanh chóng hoàn thiện bộ slide bài giảng mà không phải thực hiện thao tác tóm tắt thủ công.

\subsubsection*{Phân hệ 3: Ngân hàng câu hỏi kiểm tra đánh giá chuẩn Thang đo Bloom cải tiến}
Phân hệ thứ ba vận dụng trực tiếp thang đo nhận thức Bloom cải tiến để biên soạn bộ câu hỏi kiểm tra đánh giá. Hệ thống phân chia rõ rệt hình thức câu hỏi tương ứng với các bậc nhận thức:
\begin{itemize}
    \item Đối với năm cấp độ nhận thức đầu tiên (Nhớ, Hiểu, Áp dụng, Phân tích, Đánh giá), hệ thống khởi tạo bộ câu hỏi trắc nghiệm nhiều lựa chọn kèm đáp án và giải thích chi tiết. Các câu hỏi ở bậc áp dụng và phân tích được thiết kế dạng bài toán tình huống kỹ thuật thực tế, yêu cầu sinh viên vận dụng kiến thức lý thuyết để tính toán hoặc chỉ ra nguyên nhân sự cố.
    \item Đối với cấp độ nhận thức cao nhất (Sáng tạo), hệ thống chuyển đổi hình thức đánh giá sang các câu hỏi tự luận mở hoặc bài tập thiết kế hệ thống. Phương pháp này đòi hỏi sinh viên phải tổng hợp các tri thức đã học để đề xuất giải pháp công nghệ mới, phù hợp với yêu cầu đánh giá năng lực tư duy kỹ thuật tại Đại học Bách khoa Hà Nội.
\end{itemize}

\subsubsection*{Phân hệ 4: Hệ thống Phản biện sư phạm và Tinh chỉnh học liệu}
Phân hệ thứ tư đóng vai trò là cơ chế kiểm duyệt trực tiếp, đảm bảo quyền kiểm soát tối cao của giảng viên đối với chất lượng học liệu. Hệ thống tích hợp tác tử phản biện tự động nhằm đưa ra các câu hỏi gợi mở mang tính sư phạm để hỗ trợ giảng viên đánh giá và tối ưu hóa bài giảng (như gợi ý cân nhắc tích hợp các hoạt động thực hành ảo hoặc điều chỉnh tỷ trọng điểm số cho hoạt động nhóm). Trên cơ sở các gợi ý phản biện này, giảng viên có thể trực tiếp nhập yêu cầu điều chỉnh cụ thể vào khung tương tác và thực hiện lệnh tinh chỉnh để hệ thống tự động tái cấu trúc và tối ưu hóa nội dung học liệu theo đúng định hướng mong muốn.
